% GENERATED FILE. Edit canonical lessons, then run npm run generate:books.
% canonical-source-hash: fnv1a64:53c2cb9924e5b138
% canonical-lessons: HI-C287-nirbharkarna, HI-C287-nirbharkartahai, HI-C287-hotatoata, HI-W287-six-lines, HI-R287-depends-and-would

\chapter{It Depends, and What Would Be}
\label{ch:hi-it-depends-and-what-would-be}
\hlchaptermodality{287}

\begin{chapteropening}
\textbf{By the end of this chapter:} \emph{I can say what something depends on (… par nirbhar kartā hai), say what would happen if things were different (agar samay hotā, to maiṁ ātā), and write six connected lines.}

Complete HI-R287-depends-and-would: say what it depends on and what would have been.
\end{chapteropening}

\section[\texorpdfstring{\hi{निर्भर} \hi{करना} (nirbhar karnā)}{निर्भर करना (nirbhar karnā)}]{\hi{निर्भर} \hi{करना} (nirbhar karnā) — to depend}

\label{lesson:HI-C287-nirbharkarna}



\begin{quote}
Before the new one: \textquotedblleft{}Will we go tomorrow? If it rains, we will stay at home.\textquotedblright{} (\textbf{\hi{अगर} \hi{बारिश} \hi{होगी}, \hi{तो} \hi{हम} \hi{घर} \hi{पर} \hi{रहेंगे}}.) There is a shorter answer: \textquotedblleft{}it depends on the weather\textquotedblright{}.
\end{quote}

\subsection*{You'll want to know: \hi{निर्भर} \hi{करना}}
\textbf{\hi{निर्भर} \hi{करना}} — \emph{nirbhar karnā} — \textquotedblleft{}to depend (on)\textquotedblright{}.

\begin{quote}
\textbf{\hi{यह} \hi{मौसम} \hi{पर} \hi{निर्भर} \hi{करता} \hi{है।}}
\end{quote}

\begin{quote}
\emph{yah mausam par nirbhar kartā hai.} — \textquotedblleft{}It depends on the weather.\textquotedblright{}
\end{quote}

What it depends on takes \textbf{\hi{पर}}, \textquotedblleft{}on\textquotedblright{}. \textbf{\hi{निर्भर}} stays as it is, and \textbf{\hi{करना}} carries the ending.

\begin{itemize}
\raggedright
  \item \textbf{\hi{सब} \hi{कुछ} \hi{तुम} \hi{पर} \hi{निर्भर} \hi{करता} \hi{है।}} (\emph{sab kuch tum par nirbhar kartā hai.}) — Everything depends on you.
\end{itemize}

\begin{cousinweb}
\textbf{\hi{निर्भर}} (\emph{nirbhar}) is a tatsama. Monier-Williams gives Sanskrit \emph{nirbhara} as \textquotedblleft{}full of; excessive\textquotedblright{}; how it came to mean \textquotedblleft{}depending on\textquotedblright{} in Hindi is not recorded there. \textbf{\hi{करना}} is the everyday \textquotedblleft{}to do\textquotedblright{}, from Sanskrit \emph{karoti}, \textquotedblleft{}does\textquotedblright{} (CDIAL 2814).
\end{cousinweb}

\subsection*{Your turn}
Say these aloud:

\begin{itemize}
\raggedright
  \item \emph{nirbhar karnā}
  \item \textquotedblleft{}It depends on the weather\textquotedblright{} — \emph{yah mausam par nirbhar kartā hai}
  \item \textquotedblleft{}Everything depends on you\textquotedblright{} — \emph{sab kuch tum par nirbhar kartā hai}
\end{itemize}

\begin{tcolorbox}[breakable,colback=teal!4,colframe=teal!35!black,title={Before you move on}]
What does \textbf{\hi{निर्भर} \hi{करना}} mean? (\textquotedblleft{}To depend\textquotedblright{}.) What does the thing it depends on take? (\textbf{\hi{पर}}.)
\end{tcolorbox}
\section[\texorpdfstring{… \hi{पर} \hi{निर्भर} \hi{करता} \hi{है} (… par nirbhar …)}{… पर निर्भर करता है (… par nirbhar …)}]{… \hi{पर} \hi{निर्भर} \hi{करता} \hi{है} (… par nirbhar kartā hai) — it depends}

\label{lesson:HI-C287-nirbharkartahai}



\begin{quote}
Say \textquotedblleft{}it depends on the weather\textquotedblright{}. (\textbf{\hi{यह} \hi{मौसम} \hi{पर} \hi{निर्भर} \hi{करता} \hi{है}}.) Then \textquotedblleft{}if it rains, we will stay at home\textquotedblright{}. (\textbf{\hi{अगर} \hi{बारिश} \hi{होगी}, \hi{तो} \hi{हम} \hi{घर} \hi{पर} \hi{रहेंगे}}.)
\end{quote}

\subsection*{You'll want to know: X \hi{पर} \hi{निर्भर} \hi{करता} \hi{है}}
\begin{quote}
\textbf{\hi{हमारी} \hi{यात्रा} \hi{मौसम} \hi{पर} \hi{निर्भर} \hi{करती} \hi{है।}}
\end{quote}

\begin{quote}
\emph{hamārī yātrā mausam par nirbhar kartī hai.} — \textquotedblleft{}Our trip depends on the weather.\textquotedblright{}
\end{quote}

The subject is what depends, and \textbf{\hi{करता}} follows it: \textbf{\hi{यात्रा}} is feminine, so \textbf{\hi{करती} \hi{है}}. What it depends on takes \textbf{\hi{पर}}.

\begin{itemize}
\raggedright
  \item \textbf{\hi{यह} \hi{तुम} \hi{पर} \hi{निर्भर} \hi{करता} \hi{है।}} (\emph{yah tum par nirbhar kartā hai.}) — It's up to you.
  \item \textbf{\hi{किस} \hi{पर} \hi{निर्भर} \hi{करता} \hi{है}?} (\emph{kis par nirbhar kartā hai?}) — What does it depend on? \textbf{\hi{किस}} is the form \textbf{\hi{क्या}} and \textbf{\hi{कौन}} take before a postposition.
  \item \textbf{\hi{टिकट} \hi{का} \hi{दाम} \hi{ट्रेन} \hi{पर} \hi{निर्भर} \hi{करता} \hi{है।}} (\emph{ṭikaṭ kā dām ṭren par nirbhar kartā hai.}) — The price of a ticket depends on the train.
\end{itemize}

\subsection*{You'll want to know: it depends whether …}
\begin{quote}
\textbf{\hi{यह} \hi{इस} \hi{पर} \hi{निर्भर} \hi{करता} \hi{है} \hi{कि} \hi{बारिश} \hi{होगी} \hi{या} \hi{नहीं।}}
\end{quote}

\begin{quote}
\emph{yah is par nirbhar kartā hai ki bāriś hogī yā nahīṃ.} — \textquotedblleft{}It depends on whether it rains or not.\textquotedblright{}
\end{quote}

When what it depends on is a whole sentence, Hindi says \textbf{\hi{इस} \hi{पर}}, \textquotedblleft{}on this\textquotedblright{}, and then gives the sentence after \textbf{\hi{कि}}. It is often a question with \textbf{\hi{या} \hi{नहीं}}, \textquotedblleft{}or not\textquotedblright{}, in the future.

\begin{itemize}
\raggedright
  \item \textbf{\hi{यह} \hi{इस} \hi{पर} \hi{निर्भर} \hi{करता} \hi{है} \hi{कि} \hi{तुम} \hi{कब} \hi{आओगे।}} (\emph{yah is par nirbhar kartā hai ki tum kab āoge.}) — It depends on when you come.
\end{itemize}

\subsection*{Your turn}
Say these aloud:

\begin{itemize}
\raggedright
  \item \textquotedblleft{}It's up to you\textquotedblright{} — \emph{yah tum par nirbhar kartā hai}
  \item \textquotedblleft{}It depends on when you come\textquotedblright{} — \emph{yah is par nirbhar kartā hai ki tum kab āoge}
\end{itemize}

\begin{tcolorbox}[breakable,colback=teal!4,colframe=teal!35!black,title={Before you move on}]
\textquotedblleft{}It depends on the weather\textquotedblright{}? (\textbf{\hi{यह} \hi{मौसम} \hi{पर} \hi{निर्भर} \hi{करता} \hi{है}}.) Why \textbf{\hi{करती}} in \textbf{\hi{हमारी} \hi{यात्रा} \hi{मौसम} \hi{पर} \hi{निर्भर} \hi{करती} \hi{है}}? (\textbf{\hi{यात्रा}} is feminine.) \textquotedblleft{}It depends on whether …\textquotedblright{}? (\textbf{\hi{यह} \hi{इस} \hi{पर} \hi{निर्भर} \hi{करता} \hi{है} \hi{कि} …})
\end{tcolorbox}
\section[\texorpdfstring{\hi{अगर} \hi{समय} \hi{होता}, \hi{तो} … (agar samay hotā …)}{अगर समय होता, तो … (agar samay hotā …)}]{\hi{अगर} \hi{समय} \hi{होता}, \hi{तो} \hi{मैं} \hi{आता} (agar samay hotā, to maiṁ ātā) — what would be}

\label{lesson:HI-C287-hotatoata}



\begin{quote}
Say \textquotedblleft{}if it rains, we will stay at home\textquotedblright{}. (\textbf{\hi{अगर} \hi{बारिश} \hi{होगी}, \hi{तो} \hi{हम} \hi{घर} \hi{पर} \hi{रहेंगे}}.) Then \textquotedblleft{}used to speak\textquotedblright{}. (\textbf{\hi{बोलता} \hi{था}}.) The new form looks like that \textbf{\hi{बोलता}} with its \textbf{\hi{था}} taken away.
\end{quote}

\subsection*{You'll want to know: if it were so}
\begin{quote}
\textbf{\hi{अगर} \hi{मेरे} \hi{पास} \hi{समय} \hi{होता}, \hi{तो} \hi{मैं} \hi{आता।}}
\end{quote}

\begin{quote}
\emph{agar mere pās samay hotā, to maiṁ ātā.} — \textquotedblleft{}If I had time, I would come.\textquotedblright{}
\end{quote}

When the condition is not true, both halves take the \textbf{-\hi{ता}} form with nothing after it: \textbf{\hi{होता}}, \textbf{\hi{आता}}. It agrees as the present does: \textbf{\hi{आती}} for a woman, \textbf{\hi{आते}} for several, and \textbf{\hi{होती}} for a feminine thing.

\begin{itemize}
\raggedright
  \item \textbf{\hi{अगर} \hi{मेरे} \hi{पास} \hi{पैसे} \hi{होते}, \hi{तो} \hi{मैं} \hi{यह} \hi{घर} \hi{ख़रीदता।}} (\emph{agar mere pās paise hote, to maiṁ yah ghar kharīdtā.}) — If I had the money, I would buy this house.
  \item \textbf{\hi{अगर} \hi{मैं} \hi{तुम्हारी} \hi{जगह} \hi{होता}, \hi{तो} \hi{मैं} \hi{डाक्टर} \hi{के} \hi{पास} \hi{जाता।}} (\emph{agar maiṁ tumhārī jagah hotā, to maiṁ ḍākṭar ke pās jātā.}) — If I were you, I would go to the doctor.
\end{itemize}

There is no \textbf{\hi{ने}} with this form, even with a thing: \textbf{\hi{मैं} \hi{घर} \hi{ख़रीदता}}, never \textbf{\hi{मैंने}}.

\subsection*{You'll want to know: if it had been so}
\begin{quote}
\textbf{\hi{अगर} \hi{मुझे} \hi{पता} \hi{होता}, \hi{तो} \hi{मैं} \hi{तुम्हें} \hi{बताता।}}
\end{quote}

\begin{quote}
\emph{agar mujhe patā hotā, to maiṁ tumhẽ batātā.} — \textquotedblleft{}If I had known, I would have told you.\textquotedblright{}
\end{quote}

The same form looks back at a past that did not happen. Only the sense tells you which it is: \textquotedblleft{}if I knew\textquotedblright{} or \textquotedblleft{}if I had known\textquotedblright{}.

\begin{itemize}
\raggedright
  \item \textbf{\hi{अगर} \hi{तुम} \hi{कल} \hi{आते}, \hi{तो} \hi{अच्छा} \hi{होता।}} (\emph{agar tum kal āte, to acchā hotā.}) — If you had come yesterday, it would have been nice.
\end{itemize}

\subsection*{Your turn}
Say these aloud:

\begin{itemize}
\raggedright
  \item \textquotedblleft{}If I had the money, I would buy this house\textquotedblright{} — \emph{agar mere pās paise hote, to maiṁ yah ghar kharīdtā}
  \item \textquotedblleft{}If you had come yesterday, it would have been nice\textquotedblright{} — \emph{agar tum kal āte, to acchā hotā}
\end{itemize}

\begin{tcolorbox}[breakable,colback=teal!4,colframe=teal!35!black,title={Before you move on}]
\textquotedblleft{}If I had time\textquotedblright{}? (\textbf{\hi{अगर} \hi{मेरे} \hi{पास} \hi{समय} \hi{होता}}.) \textquotedblleft{}I would come\textquotedblright{}? (\textbf{\hi{मैं} \hi{आता}}, or \textbf{\hi{आती}}.) Does this form take \textbf{\hi{ने}}? (No: \textbf{\hi{मैं} \hi{घर} \hi{ख़रीदता}}.) \textquotedblleft{}If I had known\textquotedblright{}? (\textbf{\hi{अगर} \hi{मुझे} \hi{पता} \hi{होता}}.)
\end{tcolorbox}
\section[(six lines)]{Six lines about a trip}

\label{lesson:HI-W287-six-lines}



\begin{quote}
Your first written message gave a reader a few practised lines. This task asks for six lines, and they tell a story: a trip, to a village or anywhere else.

Pick one reader: Mira, or a friend you name. Say \textbf{\hi{ऐसा}} and \textbf{\hi{अतः}} aloud; one of them will go in.
\end{quote}

\subsection*{What to know first}
A told story is not a list of events. It has a \textbf{scene}, in \textbf{-\hi{रहा} \hi{था}} or with \textbf{\hi{जब} … \hi{तब} …}. It has \textbf{events} in the plain past, kept in order with \textbf{\hi{पहले}}, \textbf{\hi{फिर}}, \textbf{\hi{उसके} \hi{बाद}} and \textbf{\hi{आख़िर} \hi{में}}, and joined with \textbf{-\hi{कर}}. It has a \textbf{reason}, with \textbf{\hi{इसलिए}}, \textbf{\hi{क्योंकि}} or \textbf{\hi{की} \hi{वजह} \hi{से}}, and a \textbf{feeling}, with \textbf{\hi{मुझे} … \hi{लगा}} or \textbf{\hi{मैं} \hi{ख़ुश} \hi{हुआ}}.

A good one ends by looking at what might have been: \textbf{\hi{अगर} …, \hi{तो} …}, with \textbf{-\hi{ता}} in both halves.

\subsection*{Writing — connected composition}
On paper, in Devanagari, write \textbf{six lines} to your reader about a trip, real or invented, one line for each of these:

\begin{enumerate}
\raggedright
  \item the scene: when it was, with \textbf{\hi{पिछले}} (\textbf{\hi{पिछले} \hi{साल}}, \textbf{\hi{पिछले} \hi{हफ़्ते}}), and what you were doing as it began, with \textbf{-\hi{रहा} \hi{था}} or \textbf{\hi{जब} … \hi{तब} …};
  \item the first thing that happened, with \textbf{\hi{पहले}}: buying a ticket, with \textbf{\hi{टिकट}}, and when the train left, with \textbf{\hi{बजे}};
  \item what happened next, with \textbf{\hi{फिर}} or a \textbf{-\hi{कर}} form, and one thing that went wrong: \textbf{\hi{खो} \hi{गया}} or \textbf{\hi{छूट} \hi{गई}};
  \item why, with \textbf{\hi{इसलिए}}, \textbf{\hi{क्योंकि}}, \textbf{\hi{की} \hi{वजह} \hi{से}} or \textbf{\hi{वजह} \hi{यह} \hi{है} \hi{कि}};
  \item how it ended and how you felt, with \textbf{\hi{आख़िर} \hi{में}} and \textbf{\hi{मुझे} … \hi{लगा}}, \textbf{\hi{मैं} \hi{ख़ुश} \hi{हुआ}} or \textbf{\hi{मैं} \hi{उदास} \hi{हो} \hi{गया}};
  \item what might have been, with \textbf{\hi{अगर} …, \hi{तो} …} and the \textbf{-\hi{ता}} form (\textbf{\hi{अगर} \hi{मुझे} \hi{पता} \hi{होता} …}), and one word with \textbf{\hi{ऐ}} or \textbf{\hi{ः}}: \textbf{\hi{ऐसा}}, \textbf{\hi{अतः}}.
\end{enumerate}

There is no model in this lesson. Use only words you already write. The work is choosing the right past for each line, and joining the lines.

\begin{tcolorbox}[breakable,colback=teal!4,colframe=teal!35!black,title={Before you move on}]
Now read your lines as you read Mira's holiday. Mark every past verb \textbf{s} for scene and \textbf{e} for event. Line one should be mostly \textbf{s}; lines two and three mostly \textbf{e}.

Then three checks. A plain past with a thing takes \textbf{\hi{ने}} on its doer and follows the thing: \textbf{\hi{मैंने} \hi{टिकट} \hi{लिया}}, \textbf{\hi{मैंने} \hi{चाय} \hi{पी}}. A \textbf{-\hi{रहा} \hi{था}} verb never takes \textbf{\hi{ने}}. \textbf{\hi{ऐ}} has one arc above its head-line, and \textbf{\hi{ः}} has its upper dot drawn first.

Last, read line six alone. Do both halves wear \textbf{-\hi{ता}}?
\end{tcolorbox}
\section[(review)]{Review — it depends, and what would have been}

\label{lesson:HI-R287-depends-and-would}



\begin{quote}
Say \textquotedblleft{}leave early, or you'll be late\textquotedblright{}. (\emph{jaldī niklo, varnā der ho jāegī}.) Then \textquotedblleft{}of all fruits I like mangoes the most\textquotedblright{}. (\emph{mujhe phaloṁ meṁ ām sabse zyādā pasand hai}.) Then read aloud one of your six lines that has \textbf{\hi{ऐ}} or \textbf{\hi{ः}}.
\end{quote}

\subsection*{Your turn}
Say these aloud:

\begin{itemize}
\raggedright
  \item Will we go tomorrow? It depends on the weather.
  \item Is the bus cheaper than the train? It depends on the day.
  \item If I had worked all day, I would have been tired.
  \item If I had time, I would come.
  \item If it hadn't rained, we would have gone to the river.
  \item one trip that depends on something, and one that would have gone better if …
\end{itemize}

\begin{tcolorbox}[breakable,colback=teal!4,colframe=teal!35!black,title={Before you move on}]
\textquotedblleft{}It depends on\textquotedblright{}? (\textbf{… \hi{पर} \hi{निर्भर} \hi{करता} \hi{है}}.) \textquotedblleft{}It depends on whether …\textquotedblright{}? (\textbf{\hi{यह} \hi{इस} \hi{पर} \hi{निर्भर} \hi{करता} \hi{है} \hi{कि} …}) \textquotedblleft{}I would come\textquotedblright{}? (\textbf{\hi{मैं} \hi{आता}}.) \textquotedblleft{}If I had known\textquotedblright{}? (\textbf{\hi{अगर} \hi{मुझे} \hi{पता} \hi{होता}}.) Which letters did \textbf{\hi{ऐसा}} and \textbf{\hi{अतः}} bring? (\textbf{\hi{ऐ}} and \textbf{\hi{ः}}.)
\end{tcolorbox}
